\scanpage{86}

$\unsure{X|S|\text{\S}}$ schéma.

1.1. Données équivalentes

a) fibré en droites projectives \inserted{$P$} sur $X$ (fibré de Br. Sev. de rang 1) $f:P\longrightarrow X$

b) ``\deleted{\unsure{fibré de droites}} loc. libre \inserted{\unsure{rang|rg} 2 $E$} mod \deleted{\unsure{ten}}\unsure{tensoris.|torsion|tensorisation} par \unsure{module|modules} \unsure{inversible|inversibles}'' sur $X$

c) Algèbre d'Azumaya \inserted{$M$} de rang 4

\margin{si 2 inv. sur $X$}

d) Algèbre de Lie $\unsure{\mathfrak g|\mathcal G}$ forme \inserted{(fppf)} de $\mathfrak{gp}(1)_X$\\
 d') \rule{2.1cm}{0.3pt} $\mathfrak g'$ \rule{2.1cm}{0.3pt} \deleted{$\unsure{\mathfrak{pg}}$} $\unsure{\mathfrak{sgl}|\mathfrak{sl}}(1)_X$

e) forme $G$ de $\underline{\mathrm{GP}}(1)_X$\\
 \deleted{\unsure{e)}} \unsure{e$_1$)|e')} \rule{1.5cm}{0.3pt} $G'$ de $\underline{\mathrm{SL}}(2)_X$.

Si on a $P$, on en déduit $\bullet$ $E=f_*(L)$, $\underline L$ module inv. sur $P$ de degré 1 sur \unsure{ch.que|ch.} fibre ($\underline L$ défini loc. fpqc \inserted{sur $X$}, loc$^{\mathrm{lt}}$, mod tensorisation par $f^*(\lambda)$, $\lambda$ inv sur $X$)

\[ \begin{aligned}
 M&=\check E\otimes E\simeq\underline{\mathrm{End}}(E)\simeq E\otimes E\otimes(\overset2\Lambda E)^{-1}\quad\text{ext. de}\\
 &\underline O_X\;(\simeq\overset2\Lambda E\otimes(\overset2\Lambda E)^{-1})\quad\text{par }\Gamma^2E\otimes(\overset2\Lambda E)^{-1}
 \end{aligned} \]

\[ \text{ext. de }\underbrace{\mathrm{Sym}^2(E)\otimes(\overset2\Lambda E)^{-1}}_{f_*(L^{\otimes2})\otimes(\overset2\Lambda E)^{-1}\simeq f_*(\underbrace{L^{\otimes2}\otimes/\overset2\Lambda E}_{\underline t_{P/X}})}
 \text{ par }\underline O_X\;(\simeq\overset2\Lambda E\otimes(\overset2\Lambda E)^{-1}) \]

\margin{expliciter !}

\diagram{\begin{tikzcd}[ampersand replacement=\&,row sep=small]
 \mathrm{Sym}^2\check E\otimes\mathrm{Sym}^2 E\arrow[d]\\
 \mathrm{Sym}^2\check E\otimes\overset2\Lambda E\\
 \text{et}\\
 \Gamma^2\check E\otimes\Gamma^2 E\arrow[d]\\
 \Gamma^2\check E\otimes\overset2\Lambda E
 \end{tikzcd}}

\margin{\deleted{\unsure{des coefficients}} formes de Killing}

\diagram{\begin{tikzcd}[ampersand replacement=\&,column sep=small,row sep=tiny]
 \mathrm{Sym}^2\check E\otimes\mathrm{Sym}^2 E\arrow[r]\&(\overset2\Lambda\check E)^{\otimes2}\\
 \Gamma^2\check E\otimes\Gamma^2 E\arrow[ur]\&
 \end{tikzcd}}

\[ \begin{aligned}
 M&\overset{?}{\simeq}f_*(\mathcal P^1_{P/X}(\underline t_{P/X}))\quad\text{\unsure{additivement|additivem.|addition} ?}\\
 \mathfrak g&=M/\underline O_X \quad\text{avec crochet déduit du crochet }\inserted{xy-yx}\text{ dans }M\ [\underline{\mathrm{NB}}\ \underline O_X\text{ en est le centre}]\\
 &\simeq f_*(\underline t_{P/X})\quad [\text{avec crochet provenant du crochet dans }\underline t_{P/X}]
 \end{aligned} \]

\[ \begin{aligned}
 \mathfrak g'&\subset M=\unsure{E|\mathrm{tr}=0}\quad\text{sous-algèbre des éléments de trace réduite nulle}\\
 &(\mathfrak g'\longrightarrow\mathfrak g\ \text{can., iso si 2 inv. sur }X)\\
 \mathfrak g'&\subset\overset{\Vert}{\check E}\otimes E\nearrow\check E\quad\text{sous. algèbre des él. de trace nulle}
 \end{aligned} \]

\[ \begin{aligned}
 G&=\mathrm{Aut}_X P,\qquad G'=\widetilde G\quad(\text{rev. \unsure{univ.} de Chevalley})\\
 &&=\deleted{\unsure{\mathrm{Aut}}}\underline{\mathrm{SL}}(E)
 \end{aligned} \]

\margin{$\begin{aligned}G&\simeq\underline{\mathrm{Aut}}(E)/G_{mX}\\&\simeq\underline{\mathrm{Aut}}(M)\\&\simeq\underline{\mathrm{Aut}}(\mathfrak g)\\&\simeq\underline{\mathrm{Aut}}(\mathfrak g')\end{aligned}$}

$P=P(E)={}$ schéma des Idéaux à gauche qui sont ``spéciaux'', i.e. loc. fact. directs et ``non \uline{\unsure{triviaux, i.e aussi|triviaux, au sens}}'' \deleted{\unsure{non tr.}}

\[ \begin{aligned}
 &=\text{schéma des sous-Alg. de Borel de }\mathfrak g\\
 &=\rule{6cm}{0.3pt}\quad\mathfrak g'\\
 &=\rule{4cm}{0.3pt}\quad\text{Gr. de Borel de }G\\
 &=\rule{6cm}{0.3pt}\quad\text{--- }G'
 \end{aligned} \]

\margin{! attention à la Car. (notamment Car. 2 !)}

% archive sheet number: 85
